\documentclass[11pt]{article}

\usepackage[a4paper,margin=1in]{geometry}
\usepackage{amsmath,amssymb,amsthm}
\usepackage{booktabs}
\usepackage{longtable}
\usepackage{array}
\usepackage{enumitem}
\usepackage[colorlinks=true,linkcolor=blue,citecolor=blue]{hyperref}
\usepackage{titlesec}
\usepackage{setspace}
\usepackage{xcolor}

\newcommand{\note}[1]{\textcolor{blue!60!black}{\small [#1]}}
\newcommand{\code}[1]{\texttt{#1}}

\titleformat{\section}{\normalfont\Large\bfseries}{\thesection}{1em}{}
\titleformat{\subsection}{\normalfont\large\bfseries}{\thesubsection}{1em}{}

\title{\textbf{The Mathematics of the pmsimstats-ng Manuscript
Compendium}\\[0.3em]
\large A unified account of the notation, data-generating processes,
and analysis models across the eleven manuscripts under
\code{analysis/report/}}
\author{pmsimstats team}
\date{2026-08-20}

\begin{document}
\maketitle

\begin{abstract}
\noindent
The eleven manuscripts of the pmsimstats-ng compendium share one
outcome model, one family of trial designs, and one analysis-model
convention, but each paper varies a different piece of that shared
machinery to ask a different question. This document collects, in one
place and in the compendium's own notation, the mathematics that each
manuscript actually states: the response trajectory, the treatment-
exposure and carryover mechanics, the two (and, jointly, three)
architectures by which a biomarker moderates treatment response, a
latent-variable reading of the covariance-moderation architecture, the
decomposition of the treatment effect into its biological, placebo,
and natural-history parts, and the inferential machinery used to test
and calibrate the resulting interaction estimand. Where a manuscript's
own presentation is thin, inconsistent, or silent on a point the
compendium's design nonetheless raises, that gap is flagged explicitly
rather than papered over with new derivation. Every equation below is
attributed to the manuscript it came from.
\end{abstract}

\tableofcontents

%=============================================================
\section{Scope and reading order}
%=============================================================

The compendium's eleven manuscripts fall into two families. Nine share
the biomarker-treatment interaction $\beta_{bm:D}$ as their estimand
and the common outcome components of Section~\ref{sec:notation-outcome}; one
(Paper~04) targets the treatment \emph{main} effect $\beta_D$ under a
simpler, biomarker-free model; one (Paper~10) is written in
deliberately general, non-N-of-1-specific notation to isolate what
drives realized Type~I error in \emph{any} Wald test on a mixed-model
interaction coefficient. This document orders the material so that
each section can be read against the one before it:

\begin{enumerate}[label=\textbf{\arabic*.}, itemsep=2pt]
\item \textbf{Notation} (\S\ref{sec:notation}) fixes indices, the sign
  convention, and the rule that $N$ is always a total across
  randomization paths.
\item \textbf{The response trajectory} (\S\ref{sec:gompertz}, Papers
  06--07) gives the modified-Gompertz curve that every latent factor
  in the outcome follows, and the alternative curve families used to
  test whether the compendium's conclusions depend on that choice.
\item \textbf{Trial designs, exposure, and carryover}
  (\S\ref{sec:designs}, Papers 02, 05) gives the four designs, the
  treatment-exposure indicators, and the carryover decay families that
  every later section takes as fixed infrastructure.
\item \textbf{Two architectures for a biomarker-treatment interaction}
  (\S\ref{sec:architectures}, Paper 01) is the compendium's central
  methodological contribution: the same interaction can live in the
  outcome's first moment (mean moderation) or its second moment
  (covariance moderation), and the two are not the same model wearing
  different clothes.
\item \textbf{A latent-variable reading of covariance moderation}
  (\S\ref{sec:latent}, Paper 03) asks what generative mechanism a
  fitted covariance-moderation model is secretly standing in for, and
  gives the moment conditions under which a two-class mixture and a
  single MVN are observationally close.
\item \textbf{The dual-channel architecture}
  (\S\ref{sec:dualchannel}, Paper 11) activates both channels at once
  and shows their contributions to the interaction test are additive
  to first order, with the two channel weights not separately
  identifiable from a single fitted coefficient.
\item \textbf{Partitioning the treatment effect}
  (\S\ref{sec:partition}, Paper 06) decomposes the outcome itself into
  biological, placebo-belief, and natural-history components, and
  derives closed-form bias identities for what a one-component
  analysis actually estimates when the true generative process has
  three.
\item \textbf{Informative dropout} (\S\ref{sec:dropout}, Paper 09)
  gives the hazard mechanism used to censor participants as a function
  of their own symptom trajectory, and the missingness classification
  that follows from it.
\item \textbf{Inference: from a fitted coefficient to a calibrated
  test} (\S\ref{sec:inference}, Papers 08, 10, 02) collects the
  analysis-model machinery: a closed-form repeated-measures ANOVA
  derivation, a general four-channel decomposition of realized
  Type~I error, and the cluster-robust sandwich-variance family used
  to correct it.
\item \textbf{The treatment main-effect track}
  (\S\ref{sec:maineffect}, Paper 04) is the compendium's odd paper
  out: no biomarker channel at all, a linear (not Gompertz) DGP, and
  the classical closed-form power theory of Senn and of Zucker et al.
  that the rest of the compendium replaces with simulation.
\item \textbf{Shared performance measures}
  (\S\ref{sec:performance}) collects the Monte Carlo standard error,
  power, and calibration-ratio definitions used identically across
  every paper.
\end{enumerate}

Every symbol is defined at first use and, where the compendium's own
\code{NOTATION.md} governs it, cross-referenced to that governing
rule. Sign convention throughout: the three outcome components are
non-negative \emph{reductions} in symptom severity, so treatment and
interaction coefficients are negative and moderation parameters are
positive (\code{NOTATION.md} Rule 5).

%=============================================================
\section{Notation and conventions}
\label{sec:notation}
%=============================================================

\subsection{Indices and structural counts}

\begin{center}
\begin{tabular}{@{}ll@{}}
\toprule
Symbol & Meaning \\
\midrule
$i$ & participant index \\
$t$ & timepoint index (weekly measurement occasions throughout) \\
$j$, $ij$ & period or occasion index within participant (Papers 04, 06, 08) \\
$N$ & \textbf{total} number of participants, across all randomization paths \\
$P$ & number of randomization paths in a design; expected allocation $N/P$ per path \\
$n_{\mathrm{reps}}$, $n_{\mathrm{sim}}$ & Monte Carlo replicates per cell \\
$k$ & cycles per participant in a design sweep; also the Weibull shape parameter \\
\bottomrule
\end{tabular}
\end{center}

\textbf{Rule (\code{NOTATION.md} Rule 4).} $N$ is always the total
sample size across randomization paths, never a per-path count.
Table~\ref{tab:paths} gives the path count $P$ for each design; the
per-path allocation is $\lfloor N/P \rfloor$ with the remainder spread
one per path.

\begin{table}[h]
\centering
\caption{Trial designs and their randomization-path counts.}
\label{tab:paths}
\begin{tabular}{@{}lll@{}}
\toprule
Design & $P$ & Paths \\
\midrule
OL (open-label) & 1 & on drug throughout \\
CO (crossover) & 2 & drug-then-placebo, placebo-then-drug \\
OL+BDC (blinded discontinuation) & 2 & discontinuation at either of two visits \\
Hybrid & 4 & two discontinuation points $\times$ two crossover orders \\
\bottomrule
\end{tabular}
\end{table}

This distinction is not cosmetic: \code{NOTATION.md} documents a
2026-08-08 compendium-wide correction in which several drivers had
read \code{modelparam\$N} as a per-path count rather than a total.
Re-matching every design to a common total $N=70$ moved Paper 01's
Hybrid-design power from a ceiling of $\approx 0.99$ down to $0.84$
(tripling its measured carryover loss) and withdrew Paper 09's
second headline conclusion, that its most powerful design was also its
most dropout-robust --- the earlier near-immunity turned out to be a
ceiling effect at the wrong $N$. This is the compendium's clearest
demonstration that an unmatched design comparison does not merely
inflate one arm; it can manufacture an apparent robustness finding out
of nothing.

\subsection{Outcome and latent components}
\label{sec:notation-outcome}

\begin{center}
\begin{tabular}{@{}ll@{}}
\toprule
Symbol & Meaning \\
\midrule
$Y_{it}$ & observed symptom score, the modelled outcome (code column \code{Sx}) \\
$\mathrm{BL}_i$ & participant-specific baseline level \\
$BR_{it}$ & biological (pharmacological) response component \\
$PB_{it}$ & placebo-belief response component \\
$TV_{it}$ & time-variant natural-history component \\
$\varepsilon_{it}$ & observation-level residual \\
$u_i$, $b_i$ & participant random intercept \\
\bottomrule
\end{tabular}
\end{center}

\subsection{Treatment exposure and carryover}

\begin{center}
\begin{tabular}{@{}ll@{}}
\toprule
Symbol & Meaning \\
\midrule
$D_{it}$ & binary drug state: 1 on drug, 0 off drug \\
$D_{bc,it}$ & continuous exposure-decayed drug indicator (code \code{Dbc}) \\
$t_{sd}$ & time since discontinuation (code \code{tsd}) \\
$t_{od}$ & cumulative time on drug (code \code{tod}) \\
$t_{1/2}$ & carryover half-life (code \code{carryover\_t1half}) \\
$\lambda$ & carryover decay rate, $\lambda = \ln 2 / t_{1/2}$ (code \code{lambda\_cor}) \\
\bottomrule
\end{tabular}
\end{center}

\textbf{Units.} $t_{1/2}$ is quoted in \emph{weeks} in the
interaction-focused manuscripts (canonical grid $\{0, 0.5, 1.0\}$) and
in \emph{days} in the main-effect and test-procedure manuscripts, on a
pharmacokinetic scale with a 3-day baseline (1 week = 7 days).

\subsection{Biomarker, moderation, and estimand}

\begin{center}
\begin{tabular}{@{}ll@{}}
\toprule
Symbol & Meaning \\
\midrule
$B_i$ & participant-level pre-treatment biomarker (code \code{bm}) \\
$b_i$ & standardized biomarker, $(B_i - \bar B)/s_B$ (code \code{bm\_z}) \\
$c_{bm}$ & covariance-moderation parameter: on-drug $\mathrm{Cor}(B_i, BR_{it})$ \\
$\beta_{bm}$ & mean-moderation parameter: multiplier scaling the treatment effect by $b_i$ \\
$c_{bm,A}$, $c_{bm,B}$ & independent mean-channel and covariance-channel weights (dual-channel) \\
$\beta_{bm:D}$ & the target estimand, the biomarker-by-treatment interaction (code \code{bm:Dbc}) \\
$\pi$ & power, $\Pr(p < \alpha)$; $\alpha = 0.05$ throughout \\
\bottomrule
\end{tabular}
\end{center}

\textbf{Rule (\code{NOTATION.md} Rule 2).} $c_{bm}$ is a correlation;
$\beta_{bm}$ is a multiplier. They are calibrated to a common numeric
scale --- the reference value $0.45$ denotes matched moderation
strength under either architecture, because the mean-channel shift is
applied as $\beta_{bm}\, b_i\, \sigma_{BR}$ on the standardized
biomarker --- but the two symbols are not interchangeable, and a paper
using both architectures at a matched value must say so explicitly.

%=============================================================
\section{The response trajectory: modified-Gompertz progression}
\label{sec:gompertz}
%=============================================================

Every latent factor in the outcome model --- $BR$, $PB$, and $TV$ ---
is generated as a trajectory in $t$ (or in $t_{od}$, cumulative time on
drug) that rises from zero toward an asymptote. The default functional
form throughout the package is a \emph{modified Gompertz curve}; Paper
07 exists specifically to ask how much of the compendium's power and
bias findings survive if that functional-form choice is wrong.

\subsection{The modified Gompertz curve}

Paper 07's canonical statement, with a vertical offset so the curve
passes through the origin:
\begin{equation}
y(t) = \mathrm{maxr} \cdot \exp\!\left(-\mathrm{disp} \cdot
       \exp(-\mathrm{rate} \cdot t)\right), \qquad y(0) = 0
\label{eq:gompertz-07}
\end{equation}
with $\mathrm{maxr}$ the asymptotic maximum response, $\mathrm{rate}$
governing the approach speed, and $\mathrm{disp}$ the curvature near
$t=0$; the inflection occurs at $t^\ast =
\log(\mathrm{disp})/\mathrm{rate}$.

Paper 06's equivalent statement makes the offset explicit and gives it
as a normalized curve on $[0,1]$ before scaling:
\begin{equation}
C(t) = m_C \exp\!\big(-d_C \exp(-r_C t)\big) - m_C \exp(-d_C),
\label{eq:gompertz-06}
\end{equation}
for a generic component $C \in \{BR, PB, TV\}$, where $m_C$ is the
scale/asymptote parameter, $d_C$ the displacement, and $r_C$ the onset
rate; the offset $m_C\exp(-d_C)$ is precisely the value the unshifted
curve would take at $t=0$, subtracted so $C(0)=0$. Equations
\eqref{eq:gompertz-07} and \eqref{eq:gompertz-06} are the same curve;
\code{NOTATION.md}'s own shorthand is $y = \mathrm{maxr}\cdot
\exp(-\mathrm{disp}\cdot\exp(-\mathrm{rate}\cdot t))$ with a vertical
offset adjustment, matching both.

Paper 05's variant, folded into a normalized $[0,1]$ scale, is
algebraically the same curve written as a min-max normalization:
\begin{equation}
y(t) = \mathrm{maxr} \cdot
  \frac{\exp(-\mathrm{disp}\cdot\exp(-\mathrm{rate}\cdot t)) -
        \exp(-\mathrm{disp})}{1 - \exp(-\mathrm{disp})}.
\label{eq:gompertz-05}
\end{equation}

\subsection{Alternative response-curve families}

Paper 07 evaluates three alternatives, each calibrated so its
asymptotic maximum and its response level at a single anchor time
($t=5$ weeks) match the Gompertz reference:

\begin{align}
\text{Logistic:} \qquad & y(t) = \frac{\mathrm{maxr}}{1 +
  \exp(-\mathrm{rate}\cdot(t - t_0))} - \mathrm{offset}
  \label{eq:logistic} \\
\text{Hyperbolic tangent:} \qquad & y(t) = \mathrm{maxr}\cdot
  \frac{\tanh(\mathrm{rate}\cdot(t-t_0)) + 1}{2} - \mathrm{offset}
  \label{eq:tanh} \\
\text{Piecewise-linear breakpoint:} \qquad & y(t) = \mathrm{maxr}\cdot
  \min\!\left(1, \frac{t}{t_{\mathrm{bp}}}\right) +
  \mathrm{post\_slope}\cdot\max(0,\, t - t_{\mathrm{bp}})
  \label{eq:plb}
\end{align}

The logistic and hyperbolic-tangent curves are both symmetric around
an inflection at $t_0$; the piecewise-linear form is the only
non-differentiable member of the family, with a breakpoint at
$t_{\mathrm{bp}}$. \code{NOTATION.md} additionally names a
\textbf{piecewise-linear-breakpoint} family among the package's
built-in \code{br\_family} options alongside logistic and hyperbolic
tangent; Paper 07 confirms the hyperbolic-tangent form is formally a
rescaled logistic, not an independent shape family.

\subsection{How trajectory shape interacts with carryover and the
two architectures}

Under the \emph{mean-moderation} architecture the biomarker-dependent
component of the treatment effect rides the trajectory itself:
\begin{equation}
Y_{it} = \mu_0 + \beta_1 D_{it}\,(1 + \beta_{bm} B_i) + \varepsilon_{it},
\label{eq:gompertz-mean-mod}
\end{equation}
so a family with a slower early rise front-loads less of the
interaction signal onto the earliest on-drug timepoints; Paper 07
normalizes the per-timepoint moderation so its mean on-drug magnitude
equals $\sigma_{BR}$, holding the marginal interaction effect size
fixed while letting the temporal profile vary by family shape.

Under the \emph{covariance-moderation} architecture (\S\ref{sec:architectures}),
by contrast, the interaction lives in
\begin{equation}
\frac{\partial\, \mathrm{Cor}(BM, BR)}{\partial D_{bc}},
\label{eq:gompertz-cov-independence}
\end{equation}
which is independent of the trajectory family $BR$ follows: the BM--BR
covariance block is constructed directly, not derived from the mean
curve, so trajectory-shape misspecification cannot move the
covariance-channel interaction test at all. This is the structural
reason Paper 07's 16-cell factorial (four families $\times$ two
architectures $\times$ two biomarker-effect levels, $N=70$, Hybrid
OL+BDC design, $t_{1/2}=0$ to isolate the trajectory effect from
carryover) finds power differences of only a few percentage points
across families under both architectures: the mean channel is only
mildly shape-sensitive, and the covariance channel is not shape-
sensitive at all by construction. Paper 07 is explicit that this
16-cell grid deliberately fixes $t_{1/2}=0$; at $t_{1/2}>0$ the
\code{bm:Dbc}-mediated information channel would be eroded differently
by different trajectory shapes, because the off-drug exponential decay
interacts with each family's own decay profile --- a combination the
present submission does not evaluate.

%=============================================================
\section{Trial designs, treatment exposure, and carryover}
\label{sec:designs}
%=============================================================

\subsection{Simple versus complex designs}

The compendium compares four designs of increasing complexity, defined
by their randomization-path structure (Table~\ref{tab:paths}) rather
than by a closed-form exposure-sequence equation --- no manuscript in
the compendium formalizes $D_{it}$ as an explicit function of path or
sequence index; the designs are specified operationally, as
visit-by-visit exposure schedules, and $D_{it}$ is simply read off that
schedule.

\begin{description}[leftmargin=2em]
\item[OL (open-label).] Unblinded titration; no off-drug contrast, so
  it supplies no within-subject interaction information.
\item[CO (crossover).] Two-period within-subject crossover, blinded
  throughout.
\item[OL+BDC (blinded discontinuation).] Open-label titration followed
  by a blinded switch to placebo.
\item[Hybrid.] The Hendrickson (2020) reference design: open-label
  titration, blinded discontinuation, then a brief crossover; four
  randomization paths.
\end{description}

Design \emph{complexity}, in this compendium, is therefore a
statement about $P$ and about how many within-subject on/off contrasts
a design's path structure supplies --- not about a formal information
functional of $P$. Paper 05, whose remit is exactly this comparison,
contains no equation relating $P$ or switch count to the power of the
interaction estimand $\beta_{bm:D}$; its own stated target is in fact
the \emph{main} drug effect, not the interaction (that comparison is
explicitly deferred to the companion architecture manuscript, Paper
01). The one quantitative link Paper 05 does supply is between serial
correlation and the effective information in a \emph{between}-subject
contrast --- relevant to why the parallel-group RCT loses power under
AR(1) autocorrelation while the within-subject N-of-1 designs do not:
\begin{equation}
\text{effective-$N$ shrinkage factor} = \frac{1}{1 + (k-1)\rho},
\label{eq:effn-shrinkage}
\end{equation}
where $k$ is the number of repeated observations contributing to a
between-subject mean and $\rho$ is the AR(1) correlation. Because the
N-of-1 designs' primary contrast is within-subject, this shrinkage
does not act on their information the way it acts on a parallel-arm
comparator's.

\subsection{Treatment exposure: binary versus continuous}

\textbf{Rule (\code{NOTATION.md} Rule 1).} $D_{it}$ is binary;
$D_{bc,it}$ is continuous. These are different objects and must never
be substituted for one another:
\begin{equation}
D_{bc,it} = \begin{cases} 1 & D_{it} = 1 \text{ (on drug)} \\
  \exp(-\lambda\, t_{sd,it}) & D_{it} = 0 \text{ (off drug)},
\end{cases}
\label{eq:dbc}
\end{equation}
with $t_{sd,it}$ the time since discontinuation. In the limit of an
assumed zero half-life, the continuous predictor collapses exactly
onto the binary one, $D_{bc,it} \to D_{it}$ as $\hat t_{1/2} \to 0$
(Paper 02) --- the sense in which the compendium's ``Unadjusted''
analysis specification is the boundary member of the continuous-
exposure family.

\subsection{Carryover decay families}

Four decay forms are implemented, all normalized so $\phi(t_{1/2}) =
0.5$:
\begin{align}
\text{Exponential:} \qquad & \phi_{\exp}(t_{sd}) = \exp(-\lambda\,
  t_{sd}), \qquad \lambda = \frac{\ln 2}{t_{1/2}}
  \label{eq:decay-exp} \\
\text{Linear:} \qquad & \phi_{\mathrm{lin}}(t_{sd}) = \max\!\left(0,\,
  1 - \frac{t_{sd}}{2\,t_{1/2}}\right)
  \label{eq:decay-lin} \\
\text{Weibull:} \qquad & \phi_{\mathrm{weib}}(t_{sd}) =
  \exp\!\big(-(\lambda_w t_{sd})^k\big), \qquad \lambda_w =
  \frac{(\ln 2)^{1/k}}{t_{1/2}}
  \label{eq:decay-weib} \\
\text{Power:} \qquad & \phi_{\mathrm{pow}}(t_{sd}) = \max\!\left(0,\,
  \Big(1 - \frac{t_{sd}}{3\,t_{1/2}}\Big)^{\!p}\right)
  \label{eq:decay-pow}
\end{align}
The Weibull family recovers the exponential exactly at shape $k=1$; a
value $k<1$ produces a heavier elimination tail than exponential, and
$k>1$ a lighter one. Paper 02, whose remit is this decay family, tests
$k \in \{0.25, 0.5, 2.0, 4.0\}$ against the exponential default and
finds the analysis-model conclusions materially unchanged --- the
compendium's carryover-specification results are not an artifact of
assuming the exponential form. Paper 05 additionally implements a
carryover \emph{recursion} for the biological-response mean during
off-drug periods,
\begin{equation}
\mu[t] = \mathrm{base}[t] + \mu[t-1]\cdot\phi(t_{sd}),
\label{eq:carryover-recursion}
\end{equation}
so that off-drug decay compounds across consecutive off-drug
timepoints rather than resetting at each one.

\subsection{Covariance moderation with carryover}

Under the covariance-moderation architecture (\S\ref{sec:architectures}),
carryover attenuates the biomarker--response correlation itself, not
just the mean:
\begin{equation}
\mathrm{Cor}(B_i, BR_{it}) = c_{bm}\cdot\phi(t_{sd}),
\label{eq:cor-carryover}
\end{equation}
with $\phi$ any of Equations \eqref{eq:decay-exp}--\eqref{eq:decay-pow}
(Paper 02); in the exponential-only presentation of Paper 01 this is
written $\mathrm{Cor}(B_i, BR_{it}) = c_{bm}$ on drug and $c_{bm}\cdot
e^{-\lambda t_{sd}}$ off drug. The conditional-mean interpretation of
this construction, useful for connecting it to the mean-moderation
architecture, is derived in Paper 01 as
\begin{equation}
E[BR_{it} \mid B_i = b, D_{it} = 1] = \mu_{BR}(t) + c_{bm}\cdot
  \frac{\sigma_{BR}}{\sigma_B}\,(b - \mu_B),
\label{eq:cor-conditional-mean}
\end{equation}
with the off-drug analogue carrying the same decay factor
$c_{bm}\cdot\phi(t_{sd})$ in place of $c_{bm}$.

\subsection{Analysis specifications for carryover: E1--E9}

Paper 02 fixes the analysis model's carryover treatment across nine
specifications, all fit as variants of
\begin{equation}
\code{nlme::lme(Sx \textasciitilde\ bm * X + t, random = \textasciitilde 1 | ptID,
  correlation = corCAR1(form = \textasciitilde t | ptID))}
\label{eq:e-base-model}
\end{equation}
differing only in the construction of the exposure predictor $X$ and,
for the CR-variant rows, the standard-error estimator
(\S\ref{sec:sandwich}). Table~\ref{tab:especs} reproduces the full
specification table.

\begin{table}[h]
\centering\small
\caption{The nine carryover analysis specifications (Paper 02).}
\label{tab:especs}
\begin{tabular}{@{}lllc@{}}
\toprule
Code & Name & $X$ / model & Target \\
\midrule
E1 & Unadjusted & $X = D_{it}$ (binary), no carryover term & $bm{:}D_{it}$ \\
E3 & Lag-adjusted & $X = D_{it}$, plus nuisance $L_{it} = D_{i,t-1}(1-D_{it})$ & $bm{:}D_{it}$ \\
E2 & Exposure-weighted & $X = D_{bc,it}$ & $bm{:}D_{bc,it}$ \\
E7 & AIC-selected $\hat t_{1/2}$ & E2 form, $\hat t_{1/2}$ chosen by AIC over $\{0.25,0.5,1.0,2.0\}$ wk & $bm{:}D_{bc,it}$ \\
E9 & Paired-difference & two-stage OLS on per-patient on/off means & biomarker slope \\
E1cr2--E7cr2 & CR2 variants & E1/E3/E2/E7 formula, CR2 cluster-robust SE & (as above) \\
\bottomrule
\end{tabular}
\end{table}

\textbf{E3, the lag-adjusted (Jones--Kenward) specification.} The
nuisance indicator $L_{it} = D_{i,t-1}(1-D_{it})$ marks off-drug
timepoints that immediately follow an on-drug one; it enters as a main
effect only, so the interaction actually tested remains
$bm{:}D_{it}$, not $bm{:}L_{it}$.

\textbf{E7, the AIC-selected half-life.} With $M(t_{1/2})$ the E2 fit
at a candidate half-life,
\begin{equation}
\hat t_{1/2} = \operatorname*{arg\,min}_{t_{1/2}\in\mathcal T}
  \mathrm{AIC}\big(M(t_{1/2})\big), \qquad \mathrm{AIC}(M) =
  -2\log L(M) + 2p, \qquad \mathcal T = \{0.25, 0.5, 1.0, 2.0\},
\label{eq:e7-aic}
\end{equation}
with the reported coefficient, SE, and $p$-value read from
$M(\hat t_{1/2})$ alone --- a post-selection-inference risk Paper 02
notes explicitly (Berk et al., 2013): treating $\hat t_{1/2}$ as fixed
in advance understates the true sampling variability of the reported
interaction estimate.

\textbf{E9, the paired-difference specification.} A correlation-
structure-agnostic two-stage estimator. For patient $i$ with
$n_i^{\mathrm{on}}, n_i^{\mathrm{off}} > 0$,
\begin{equation}
\bar S_{x,i}^{\mathrm{on}} = \frac{1}{n_i^{\mathrm{on}}}
  \sum_{t:\,D_{it}=1} S_{x,it}, \qquad
\bar S_{x,i}^{\mathrm{off}} = \frac{1}{n_i^{\mathrm{off}}}
  \sum_{t:\,D_{it}=0} S_{x,it}, \qquad
\Delta_i = \bar S_{x,i}^{\mathrm{on}} - \bar S_{x,i}^{\mathrm{off}},
\label{eq:e9-diffs}
\end{equation}
followed by a second-stage OLS regression on the biomarker,
\begin{equation}
\Delta_i = \beta_0 + \beta_1\, bm_i + \varepsilon_i, \qquad i =
  1,\dots,N',
\label{eq:e9-ols}
\end{equation}
with the interaction test on $\beta_1$ and $N' \le N$ patients
contributing (those without both on- and off-drug observations are
dropped).

An earlier, retired pair of specifications, E5 (adding
$bm{:}L_{it}$) and E6 (adding $bm{:}t_{sd}$ as a direct
carryover-by-biomarker interaction), underperformed the Unadjusted
specification at every cell tested and were replaced by E7/E9; they
are recorded here as historical context, not part of the active
nine-specification set.

%=============================================================
\section{Two architectures for a biomarker-treatment interaction}
\label{sec:architectures}
%=============================================================

Paper 01 is the compendium's central methodological result: the same
target estimand, $\beta_{bm:D}$, can be encoded into the data-
generating process in two structurally different ways, and the two
encodings are \emph{not} numerically interchangeable once carryover is
present.

\subsection{Mean moderation (Architecture A)}

The biomarker scales the treatment effect directly, in the outcome's
conditional mean:
\begin{equation}
Y_{it} = \mu_0 + \beta_1\, D_{it}\,(1 + \beta_{bm} B_i) +
  \varepsilon_{it}.
\label{eq:archA}
\end{equation}
This is algebraically equivalent to the standard additive
effect-modification form of Kent et al.,
\begin{equation}
Y_{it} = \beta_0 + \beta_1 D_{it} + \beta_2 B_i + \beta_3 D_{it}
  B_i + \varepsilon_{it}, \qquad \beta_3 = \beta_1\,\beta_{bm},
\label{eq:archA-kent}
\end{equation}
identifying the interaction coefficient $\beta_3$ with the product of
the base treatment effect and the moderation multiplier. The mean-
difference-between-participants identity that follows,
\begin{equation}
E[Y \mid b_1, D] - E[Y \mid b_2, D] = \beta_1\,\beta_{bm}\,(b_1 -
  b_2)\, D,
\label{eq:archA-contrast}
\end{equation}
shows that under mean moderation, carryover (which shrinks $D$'s
continuous analogue $D_{bc}$ toward zero) shrinks the moderation
contrast \emph{proportionally} for every biomarker level --- the
mechanism is not selectively erased for any particular pair of
participants.

\subsection{Covariance moderation (Architecture B)}

The biomarker and the biological-response component are instead drawn
jointly from a bivariate normal whose correlation, not whose mean, is
treatment-state dependent:
\begin{equation}
\begin{pmatrix} B_i \\ BR_{it} \end{pmatrix} \sim
\mathrm{MVN}\!\left(
  \begin{pmatrix}\mu_B \\ \mu_{BR}(t, D_{it})\end{pmatrix},\,
  \Sigma(D_{it})\right), \qquad
\mathrm{Cor}(B_i, BR_{it}) = \begin{cases}
  c_{bm} & D_{it}=1 \\
  c_{bm}\cdot e^{-\lambda t_{sd}} & D_{it}=0.
\end{cases}
\label{eq:archB}
\end{equation}
The corresponding conditional-mean expressions on and off drug are
Equations \eqref{eq:cor-conditional-mean} and its carryover-decayed
analogue,
\begin{equation}
E[BR \mid b, D=0, t_{sd}] = \mu_{BR,\mathrm{off}}(t_{sd}) +
  c_{bm}\cdot e^{-\lambda t_{sd}}\cdot\frac{\sigma_{BR}}{\sigma_B}\,
  (b - \mu_B).
\label{eq:archB-offdrug}
\end{equation}
Here the biomarker-dependent part of the conditional mean decays
\emph{exponentially} with time off drug, not proportionally: as
$t_{sd}$ grows, the differential signal between biomarker strata is
eliminated entirely, in a way Equation~\eqref{eq:archA-contrast}'s
mean-moderation contrast never is. This structural asymmetry is the
reason the two architectures diverge sharply once carryover enters:
in Paper 01's matched-$N=70$ headline result, carryover power loss
from $t_{1/2}=0$ to $t_{1/2}=1.0$ weeks in the Hybrid design is
1.4--6.9 percentage points under mean moderation but 5.1--15.4 points
under covariance moderation, with a difference-of-differences
$z=2.80$ (Hybrid) and $z=7.64$ (OL+BDC).

\subsection{The common analysis model}

Both architectures are fit identically:
\begin{equation}
\code{Sx \textasciitilde\ bm + t + Dbc + bm:Dbc}
\label{eq:common-analysis}
\end{equation}
via \code{nlme::lme} with \code{corCAR1(form = \textasciitilde
t|ptID)} residual correlation and a participant random intercept ---
the shared inferential machinery of \S\ref{sec:inference} does not
know, and cannot distinguish, which architecture generated the data
it is fitting.

\note{Gap closed.} Paper 01's Methods now state the closed-form
reference value explicitly, $\theta_{\mathrm{true}} = -c_{bm}\cdot
\sigma_{BR}$, alongside its bias performance measure, with a citation
to the same calibration used in Paper 02
(Equation~\eqref{eq:theta-true} below) and Paper 06
(Equation~\eqref{eq:beta-target}). This is consistent with Paper 01's
reported mean estimates of $\hat\beta_{bm:D_{bc}} \approx -0.23$ to
$-0.28$; the manuscript itself does not evaluate $\sigma_{BR}$
numerically, so the identity is stated symbolically there, exactly as
in the two companion papers.

%=============================================================
\section{A latent-variable reading of covariance moderation}
\label{sec:latent}
%=============================================================

Paper 03 asks what mechanistic story a fitted covariance-moderation
model could be a moment-matched summary of, if the true underlying
process is a discrete responder/non-responder dichotomy rather than a
continuously varying correlation.

\subsection{The faithful latent-class data-generating process}

If a candidate predictive biomarker indexes an unobserved
responder/non-responder class,
\begin{equation}
Z_i \mid B_i \sim \mathrm{Bernoulli}\big(\pi(B_i)\big), \qquad
  BR_{it} \mid Z_i = z \sim f_z(t, D_{bc,it}),
\label{eq:latent-class}
\end{equation}
with $Z_i$ the unobserved class label and $\pi(B_i) = \Pr(Z_i=1\mid
B_i)$ the class-membership probability, now given an explicit closed
form in Paper 03's Notation section,
\begin{equation}
\pi(B_i) = \mathrm{expit}(\psi_0 + \psi_g\, b_i), \qquad b_i =
  \frac{B_i - \bar B}{s_B},
\label{eq:gating-logistic}
\end{equation}
with $\psi_g$ the \textbf{gating slope} (called \code{gs} in the
code, and now given the symbol $\psi_g$ in the manuscript rather than
left unnamed) and $\psi_0$ fixed so $\pi(0)=0.5$ throughout. $f_0,
f_1$ are class-conditional response distributions differing in
drug-response parameters. Within-class temporal structure follows the
same modified-Gompertz form of \S\ref{sec:gompertz}, with class-
specific parameters substituted.

\note{Notation caveat, resolved.} Paper 03 previously referred to
gating-function steepness only informally, as \code{gs} in prose and
code, with no manuscript symbol attached and no closed form for
$\pi(B_i)$ ever written out --- a gap this compendium originally
flagged because a symbol $\gamma$ \emph{does} appear elsewhere in the
compendium (Paper 06, \S\ref{sec:partition}) as the unrelated $BR
\times PB$ subadditivity coefficient, inviting confusion between the
two. Paper 03 now introduces $\psi_g$ and
Equation~\eqref{eq:gating-logistic} explicitly, and its Notation
section states in so many words that $\psi_g$ is not to be confused
with Paper 06's $\gamma$. \code{NOTATION.md} has been updated to match:
$\gamma$ is now documented compendium-wide as Paper 06's coefficient
only, and $\psi_g$ is added as its own entry for Paper 03's gating
slope.

\subsection{The MVN approximation to mixture-induced moderation}

Under within-class independence between $B_i$ and $BR_{it}$, a two-
component mixture with class-dependent means and common within-class
covariance produces a marginal joint distribution of $(B_i, BR_{it})$
whose first two moments are reproduced, to leading order, by a single
MVN with $\mathrm{Cor}(B, BR) = c_{bm}$, where
\begin{equation}
c_{bm}^2 = \underbrace{\frac{\pi(1-\pi)(\mu_B^1 - \mu_B^0)^2}
  {\mathrm{Var}(B)}}_{\text{between-class fraction in } B} \; \cdot \;
  \underbrace{\frac{\pi(1-\pi)(\mu_{BR}^1 - \mu_{BR}^0)^2}
  {\mathrm{Var}(BR)}}_{\text{between-class fraction in } BR}.
\label{eq:latent-mvn-approx}
\end{equation}
That is, $c_{bm}^2$ is the \emph{product} of the two components'
between-class variance fractions, so the marginal correlation an MVN
can reproduce is bounded above by the geometric mean of the two
intra-class correlations. Equivalently: Paper 01's Architecture B is,
under this reading, the leading-order MVN signature of a latent-class
mixture whose class-conditional means separate on drug and whose
separation is eroded by carryover off drug --- not necessarily a
directly generative covariance mechanism in its own right. This
reinterpretation gives covariance moderation a mechanistic story
(discrete class membership) in addition to its literal one (a
continuously specified correlation), and it identifies the conditions
under which the approximation breaks down: strongly bimodal responses,
a step-function (rather than smooth) $\pi(B_i)$, or class-varying
(rather than common) within-class covariance.

\subsection{Estimation and its two competing test statistics}

The class-\emph{unaware} analysis is the compendium's common model,
Equation~\eqref{eq:common-analysis}. The class-\emph{aware} analysis
fits a two-class growth-mixture model (\code{lcmm::hlme}) and extracts
two candidate signals: a Wald test on the biomarker's coefficient in
the class-gating equation, and a within-class Wald test on
$bm{:}D_{bc}$. A model-comparison gate is applied before either is
trusted:
\begin{equation}
\text{prefer } ng=2 \text{ over } ng=1 \iff \Delta\mathrm{BIC} > 0
  \ \ (\text{or, more conservatively,} > 6),
\label{eq:bic-gate}
\end{equation}
following Raftery (1995), with the gating test interpreted only
conditional on that selection. In Paper 03's 240-replicate pilot at
$N=70$, the unconditional class-gating Wald test is markedly
anti-conservative under the null (Type~I error $0.197$, 95\% CI
$[0.145, 0.250]$), while the BIC-conditional version is
under-powered relative to the class-unaware baseline (power $0.025$ at
gating slope $1.5$ versus $0.542$ for the ordinary $bm{:}Dbc$ Wald
test on the same cell) --- the class-aware machinery, in this pilot,
buys correct calibration at a substantial power cost, and is not
obviously the better default test.

%=============================================================
\section{The dual-channel architecture}
\label{sec:dualchannel}
%=============================================================

Paper 11 activates the mean and covariance channels of
\S\ref{sec:architectures} simultaneously, with independent weights
$c_{bm,A}$ (mean channel) and $c_{bm,B}$ (covariance channel).

\subsection{Construction}

The covariance channel is built into the MVN sigma matrix exactly as
in Equation~\eqref{eq:archB}, with $c_{bm}$ replaced by $c_{bm,B}$.
The mean channel is then applied as a post-draw additive shift at
on-drug timepoints,
\begin{equation}
\Delta BR_{it} = c_{bm,A}\cdot z_{B,i}\cdot\sigma_{BR}\cdot
  \mathbf{1}[\text{on-drug}], \qquad z_{B,i} =
  \frac{B_i - \mu_B}{\sigma_B}.
\label{eq:dualchannel-shift}
\end{equation}

\subsection{The channels add on the covariance scale}

Each channel's on-drug contribution to $\mathrm{Cov}(B, BR)$ can be
computed separately:
\begin{align}
\text{Covariance channel alone:} \qquad &
  \mathrm{Cov}(B, BR \mid \text{on-drug}) = c_{bm,B}\,\sigma_B\sigma_{BR}
  \label{eq:dc-covB} \\
\text{Mean channel alone:} \qquad &
  \mathrm{Cov}(B, \Delta BR) = c_{bm,A}\,\sigma_B\sigma_{BR}
  \label{eq:dc-covA}
\end{align}
and, because the combined-architecture on-drug response is the sum of
the correlated draw and the mean shift, the two contributions add:
\begin{equation}
\mathrm{Cov}(B, BR_C \mid \text{on-drug}) = (c_{bm,A} +
  c_{bm,B})\,\sigma_B\sigma_{BR}.
\label{eq:dc-covC}
\end{equation}
To first order the interaction test's noncentrality is likewise
additive across channels, $\eta_C \approx \eta_A + \eta_B$, and Paper
11 reports empirical \emph{super}-additivity of power itself relative
to the natural probability-scale benchmark: for rejection probability
$\pi(\eta)$ as a function of noncentrality,
\begin{equation}
\pi(\eta_A + \eta_B) - \pi(0) \; > \;
  \big[\pi(\eta_A) - \pi(0)\big] + \big[\pi(\eta_B) - \pi(0)\big].
\label{eq:dc-superadditivity}
\end{equation}
At $(c_{bm,A}, c_{bm,B}) = (0.22, 0.22)$ in the Hybrid design, at the
matched total $N=70$, for example, observed power is $0.577$ against
an additive benchmark $p(0.22,0)+p(0,0.22)-p(0,0) = 0.214+0.138-0.040
= 0.312$, an excess of $0.265$ --- direct confirmation of
Equation~\eqref{eq:dc-superadditivity}.

\subsection{Boundary reductions and an identifiability consequence,
now stated}

Setting either weight to zero recovers the corresponding pure
architecture exactly, by direct substitution into
Equations~\eqref{eq:dualchannel-shift} and~\eqref{eq:archB}:
$c_{bm,B}=0$ recovers mean moderation; $c_{bm,A}=0$ recovers
covariance moderation. Paper 11 treats these boundaries as internal
validation gates --- a correctly implemented combined driver must
reproduce the companion manuscripts' pure-architecture power values
there, run under matched design, sample-size, and calibration
settings. Equation~\eqref{eq:dc-covC} shows that the combined on-drug
covariance --- and hence the fitted $bm{:}D_{bc}$ coefficient ---
depends on $c_{bm,A}$ and $c_{bm,B}$ only through their \emph{sum}, so
a single fitted interaction coefficient cannot separately recover the
two channel weights; this non-identifiability, previously left
unstated in Paper 11's \S2.2 (``Boundary reductions and
identifiability''), is now given as an explicit consequence in the
manuscript's \S3, immediately after the additive-covariance
derivation.

\note{A second gap, investigated rather than left flagged.} Paper
11's driver (\code{analysis/scripts/dgp-combined/
01-run-combined-factorial.R}) already allocates $N$ as a total, via
the same \code{generate\_data\_multi\_path()} helper the companion
paper uses, and its stored output already includes the matched
$N=70$ cells --- so the sample-size-convention explanation an earlier
version of \S4.2 gave for the boundary mismatch does not hold up.
Reading the $N=70$, E2-specification boundary cells directly from
that stored output and comparing them to Paper 01's $N=70$ baselines
shows the two still do \emph{not} agree: at $t_{1/2}=0$ the
mean-moderation boundary reads $0.595$/$0.738$/$0.612$ (CO/Hybrid/
OL+BDC) against Paper 01's $0.739$/$0.842$/$0.751$, a gap of 10--19
percentage points, far beyond the $\approx0.013$--$0.016$ MCSE at
these cell sizes. The manuscript's \S4.2 has been rewritten to report
this finding rather than the superseded $N$-convention explanation,
and now names the likely locus of the mismatch (a calibration
difference in \texttt{resp\_param}, \texttt{baseline\_param}, or the
nuisance correlations \texttt{c.cf1t}/\texttt{c.cfct}/\texttt{rho}
between the two drivers) as an open item for a follow-up revision.
Paper 11's own $3\times3$ channel-weight power table has likewise been
recomputed on the correct $N=70$ data already sitting in that same
output file, in place of the earlier, uncorrected $N=35$ table.

%=============================================================
\section{Partitioning the treatment effect}
\label{sec:partition}
%=============================================================

Paper 06 asks what a single-indicator (``lumped'') analysis of the
treatment effect and its biomarker interaction actually estimates when
the true generative process has three distinguishable components.

\subsection{The outcome decomposition}

\begin{equation}
Y_{it} = \mathrm{BL}_i - \big[BR_{it} + PB_{it} + TV_{it}\big] +
  \varepsilon_{it},
\label{eq:decomposition}
\end{equation}
with each component following the modified-Gompertz form of
Equation~\eqref{eq:gompertz-06}, and $PB_{it}$ additionally scaled by
an \emph{expectancy multiplier} $\eta(\mathrm{phase})$ that
operationalizes belief-driven placebo response (Kirsch, 1985):
\begin{equation}
PB(t \mid \mathrm{phase}) = \eta(\mathrm{phase})\cdot m_{PB}
  \exp\!\big(-d_{PB}\exp(-r_{PB}t)\big) + \text{(offset)}, \qquad
\eta = \begin{cases} 1.0 & \text{open-label, on drug} \\
  \approx 0.5 & \text{blinded discontinuation} \\
  \in(0,0.5) & \text{open-label placebo.}
\end{cases}
\label{eq:eta-scaling}
\end{equation}
Because $\eta$ scales both the mean and SD of $PB$ by the same factor,
$\mathrm{SD}(PB\mid\mathrm{phase}) = \eta(\mathrm{phase})\cdot
\sigma_{PB}^{\mathrm{baseline}}$, only the \emph{product} $\eta\cdot
m_{PB}$ is identified from a two-phase (open-label + blinded) design;
recovering $m_{PB}$ itself requires an open-label-placebo phase or a
direct expectancy measurement.

A generalized, non-additive form allows a $BR\times PB$ interaction:
\begin{equation}
Y_{it} = \mathrm{BL}_i - \big[BR_{it} + PB_{it} + TV_{it} +
  \gamma\cdot BR_{it}\cdot PB_{it}\big] + \varepsilon_{it},
\label{eq:decomposition-gamma}
\end{equation}
with \emph{subadditivity} of the two response mechanisms defined by
the sign condition $\partial Y/\partial(BR\cdot PB) > 0$ on the
reduction scale. Paper 06 explicitly distinguishes true $\gamma\ne0$
subadditivity from an \emph{apparent} subadditivity that a latent
responder-class gating function $\pi(B)$ (\S\ref{sec:latent}) could
generate through confounding, without any real $BR$--$PB$ interaction
being present.

\subsection{What a one-component analysis actually estimates}

Writing the within-participant symptom reduction as $R_{it} =
\mathrm{BL}_i - Y_{it} = BR_{it} + PB_{it} + TV_{it} - \varepsilon_{it}$
and fitting the one-component model $R_{it} = \alpha + \beta_D D_{it} +
g(t) + u_i + e_{it}$, Paper 06 derives:
\begin{equation}
\operatorname*{plim}\hat\beta_D = m_{BR} + b_{PB} + b_{TV},
\label{eq:plim-betaD}
\end{equation}
where $b_{PB} = E[PB\mid\text{on}] - E[PB\mid\text{off}]$ and $b_{TV}$
is the natural-history contribution that fails to cancel under the
contrast used (zero for a clean on/off within-participant contrast;
$E[TV]$ itself for a baseline-anchored pre--post analysis). Both terms
are generally non-negative for an improving outcome, so the
one-component estimate systematically \emph{over-states} the
pharmacological effect $m_{BR}$ --- and this is a property of the
estimator's probability limit, unaffected as $N\to\infty$, not a
finite-sample artifact correctable by more data.

The same logic applied to the biomarker slope gives the compendium's
key attenuation/inflation identity:
\begin{equation}
\beta_{bm}^{\mathrm{lumped}} = \beta_{bm}^{BR} + w_{PB}\,\beta_{bm}^{PB}
  + w_{TV}\,\beta_{bm}^{TV}, \qquad
\beta_{bm}^{C} = \frac{\operatorname{Cov}(m_C, B)}{\operatorname{Var}(B)}
  \ \ (C \in \{BR, PB, TV\}),
\label{eq:lumped-slope}
\end{equation}
with $w_{PB}, w_{TV}\ge0$ design-determined loadings. A biomarker
correlated with placebo responsiveness or with the natural-history
trajectory therefore registers a non-zero lumped slope even when its
true pharmacological association $\beta_{bm}^{BR}$ is exactly zero ---
and, because Equation~\eqref{eq:lumped-slope} is an \emph{identity},
no amount of additional sample size can separate a genuinely
pharmacological biomarker signal from a confounded one whenever
$w_{PB}\beta_{bm}^{PB} + w_{TV}\beta_{bm}^{TV}\ne0$. Under
$\gamma\ne0$ subadditivity (Equation~\eqref{eq:decomposition-gamma}),
a further first-order correction proportional to
$\operatorname{Cov}(m_{BR}\cdot m_{PB}, B)$ enters the slope; the
identity above is the leading-order behavior in the additive case.

The operative confound is a biomarker's correlation with individual-
level \emph{variance} in placebo responsiveness, not with its
\emph{mean}:
\begin{equation}
\text{contamination driver} = c_{bm,PB}\cdot\sigma_{PB}, \qquad
  \text{not } c_{bm,PB}\cdot m_{PB}.
\label{eq:contamination-driver}
\end{equation}
Paper 06's calibration target throughout its production studies is
\begin{equation}
\beta_{bm:D_{bc}} = -\frac{c_{bm}\,\sigma_{BR}}{\sigma_{bm}} = -2.25,
\label{eq:beta-target}
\end{equation}
against which Study A (uncontaminated biomarker, $c_{bm,PB}=0$) shows
bias within one Monte Carlo standard error of zero at every sample
size tested ($N \in \{35,70,200,300\}$), confirming
Equation~\eqref{eq:lumped-slope} collapses to zero bias exactly when
$\beta_{bm}^{PB} = \beta_{bm}^{TV} = 0$; a contaminated-biomarker pilot
at $c_{bm,PB} \in \{0, 0.15, 0.30, 0.45\}$ (constrained below $0.45$
by positive-definiteness of the resulting covariance matrix) shows
bias growing from $-0.06$ to $+0.51$ (about 30 MCSEs) as contamination
increases, with the sign flip explained by an \emph{attenuation-
toward-zero} mechanism: high-$c_{bm,PB}$ participants' elevated
placebo response is partially absorbed by the biomarker main effect,
leaving the interaction coefficient itself pulled toward zero rather
than pushed away from it.

Critically, Paper 06 emphasizes throughout that this bias lives in the
\emph{location} of $\hat\beta_{bm:D_{bc}}$, not in the test's power or
rejection rate --- a significance test of the lumped coefficient
cannot, even in principle, detect this bias, because the bias does not
move the point estimate toward the null.

\subsection{Recovery, and its dependence on trial design}

A phase-augmented analysis model is the smallest step toward
component-aware inference:
\begin{equation}
\code{Sx \textasciitilde\ bm + t + Dbc + phase + Dbc:phase + bm:Dbc +
  bm:Dbc:phase}
\label{eq:phase-augmented}
\end{equation}
with random intercept and \code{corCAR1} residual correlation, where
\code{phase} (open-label / blinded / post-crossover) absorbs
systematic expectancy differences and \code{Dbc:phase},
\code{bm:Dbc:phase} test whether the apparent drug effect and its
moderation shrink under blinding. Whether this --- or full component-
matched Gompertz-basis fitting --- actually recovers the true
pharmacological slope $\beta_{bm}^{BR}$ is shown to be
\emph{design-dependent}, not a property of decomposition per se: under
the compendium's standard Hybrid design, belief-covariate
decomposition, blinded-stratum contrasts, and full Gompertz-basis
decomposition all fail to recover $\beta_{bm}^{BR}$, performing no
better (and sometimes worse) than the lumped analysis, because
$D_{bc}$ and the expectancy multiplier $\eta$ are positively
correlated ($\approx+0.6$) under that design. A belief-decoupling
\emph{balanced-placebo} design --- three open-label, six covert
on-drug ($\eta=0$), and six open-placebo ($\eta=1$) timepoints ---
reduces that correlation to $\approx-0.66$ and restores recovery: at
matched $N$, decomposition-analysis bias stays within one MCSE of zero
across $c_{bm,PB}\in\{0,0.1,0.2,0.3\}$ while the lumped estimator's
bias grows monotonically over the same grid.

%=============================================================
\section{Informative dropout}
\label{sec:dropout}
%=============================================================

Paper 09 evaluates a censoring mechanism whose hazard depends on the
participant's own symptom trajectory, crossed with the four trial
designs of \S\ref{sec:designs}.

\subsection{The dropout hazard}

Marginal per-timepoint censoring probability is the union of two
independent components:
\begin{equation}
\Pr(\mathrm{drop}_{it}) = 1 - \big(1 - p^{\mathrm{flat}}_{it}\big)
  \big(1 - p^{\mathrm{bias}}_{it}\big),
\label{eq:dropout-hazard}
\end{equation}
with $p^{\mathrm{flat}}_{it}\propto t$ depending on elapsed time alone
(calibrated to a marginal dropout fraction of $\beta_0(1-\beta_1)$)
and $p^{\mathrm{bias}}_{it}\propto t\,(\Delta_{Sx}(t)+100)^{e_{b1}}$
depending on the symptom change $\Delta_{Sx}(t)$ since baseline
(calibrated to $\beta_0\beta_1$), where $\beta_0$ is the overall
dropout rate, $\beta_1$ the fraction attributable to the symptom-
biased mechanism, and $e_{b1}=2$ throughout. The additive shift of
$100$ inside the biased component keeps its convex argument positive
across the realized range of symptom scores, so that large
improvement lowers dropout probability while worsening (or a small
change) raises it. Once a participant is censored at a timepoint,
every later timepoint for that participant is also set missing
(monotone/absorbing dropout).

\subsection{Missingness classification}

Because the flat component depends only on a fully observed design
variable (elapsed time), it induces missingness that is ignorable for
likelihood-based inference on its own. Because the biased component's
censoring indicator is a function of $\Delta_{Sx}(t)$ --- the same
quantity realized at, and thereby rendered missing by, that
timepoint --- Paper 09 argues the mechanism is MNAR for every dropout
pattern with $\beta_1 > 0$: missingness of $Y_{it}$ depends on the
value of $Y_{it}$ itself, the defining condition for missing-not-at-
random data. No manuscript in the compendium gives a formal
probability statement for MCAR/MAR/MNAR beyond this argument; the
classification rests on the hazard construction of
Equation~\eqref{eq:dropout-hazard} alone.

\subsection{Interaction with design}

The estimand is the same interaction test as elsewhere,
$\pi = \Pr(p<0.05)$ for $H_0: \beta_{bm:D}=0$, fit post-dropout. Paper
09 documents a ``happy accident'' selection effect, described
qualitatively rather than by a closed-form path-conditional formula:
participants with unambiguous open-label response are the most
readily resolved and most prone to drop out early, which enriches the
later, randomized (crossover/discontinuation) blocks for participants
whose open-label response was ambiguous --- precisely the subgroup
whose crossover contrasts are most informative for the interaction
test. This selection effect is design-dependent: under the OL design,
where $D_{bc}$ is identically one throughout, there is no within-
subject contrast for $bm{:}D_{bc}$ at all, so dropout there can only
reduce power, never selectively concentrate information. A
quantified, per-path decomposition of the happy-accident effect is
explicitly scoped for a future production-grade rerun and is not
delivered in the version inspected here; as with several other design
comparisons in this compendium (\S\ref{sec:notation}), an earlier
mismatched-$N$ run of this same comparison had manufactured an
apparent robustness finding (the Hybrid design appearing dropout-
immune) that a matched-$N=70$ re-run showed was a ceiling effect
rather than a real property of the design.

%=============================================================
\section{Inference: from a fitted coefficient to a calibrated test}
\label{sec:inference}
%=============================================================

\subsection{A closed-form comparator: repeated-measures ANOVA}
\label{sec:rmanova}

Paper 08's principal analytic contribution, stated for comparison
against the compendium's simulation-based test procedures, is a full
closed-form treatment of the classical split-plot RM-ANOVA under a
dichotomized biomarker. The model:
\begin{equation}
Y_{ijlk} = \mu + \alpha_i + \beta_j + (\alpha\beta)_{ij} +
  \pi_{l(i)} + \tau_k + (\alpha\tau)_{ik} + (\beta\tau)_{jk} +
  (\alpha\beta\tau)_{ijk} + e_{ijlk},
\label{eq:rmanova-model}
\end{equation}
for biomarker stratum $i$, treatment phase $j$, subject $l$, and
timepoint $k$, with $\pi_{l(i)}\sim N(0,\sigma_\pi^2)$ the random
subject effect and $e_{ijlk}\sim N(0,\sigma_e^2)$ the within-subject
error under compound-symmetric sphericity. The biomarker--treatment
interaction $(\alpha\beta)_{ij}$ is the target; its test statistic is
\begin{equation}
F_{AB} = \frac{MS_{AB}}{MS_{B\times S(A)}} \sim F(df_{AB},\,
  df_{B\times S(A)}) \text{ under } H_0,
\label{eq:rmanova-f}
\end{equation}
which specializes, for $a=b=2$ strata/phases and $T$ within-phase
timepoints, to a fully closed form. Writing the interaction contrast
as $\hat L = (\bar\Delta_1 - \bar\Delta_2)$ with $\bar\Delta_i$ the
stratum-$i$ mean of each subject's own phase-1-minus-phase-2 within-
subject difference,
\begin{equation}
F_{AB} = \frac{nT\,\hat L^2}{4\,\hat\sigma_e^2}, \qquad t_{AB} =
  \sqrt{\tfrac{nT}{4}}\cdot\frac{\hat L}{\hat\sigma_e}, \qquad
  t_{AB}^2 = F_{AB}, \qquad df_1 = 1,\ df_2 = 2(n-1),
\label{eq:rmanova-f-22}
\end{equation}
recovering the algebraic equivalence of the $2\times2$ interaction
$F$-test to a two-sample $t$-test on within-subject differences. Power
follows the noncentral $F$ distribution with noncentrality
\begin{equation}
\lambda = \frac{nT\,\Delta_0^2}{4\sigma_e^2}, \qquad 1-\beta =
  \Pr\!\big[F_{AB} > F_{1,\,2(n-1),\,1-\alpha} \mid \lambda\big],
\label{eq:rmanova-power}
\end{equation}
where $\Delta_0$ denotes the true value of the contrast $\hat L$
under $H_1$, an effect size on the raw outcome scale. Paper 08
originally called this quantity $\delta$, which collided with the
unrelated \emph{standardized bias} $\delta$ of the general calibration
framework in \S\ref{sec:calibration} below; the manuscript has been
relabeled to $\Delta_0$, with an explicit note distinguishing it from
Paper 10's $\delta$, so the symbol $\delta$ is now reserved
compendium-wide for standardized bias alone. Paper 08 finds this
closed form is the compendium's only fully analytic power result;
GEE, Mancl--DeRouen, and the LME/\code{corCAR1} model it compares
against are evaluated purely by simulation, with the sandwich/CR2
machinery of \S\ref{sec:sandwich} providing the small-sample
correction.

\subsection{A general decomposition of realized Type~I error}
\label{sec:calibration}

Paper 10, written in deliberately general Wald-test notation
(applicable to any fixed-effect interaction coefficient, not only
$\beta_{bm:D}$), decomposes an estimator into a standardized component:
\begin{equation}
\hat\theta = \mu + \sigma Z^\ast, \qquad \mu = E(\hat\theta), \qquad
  \sigma = \mathrm{SD}(\hat\theta), \qquad Z^\ast =
  \frac{\hat\theta - \mu}{\sigma},
\label{eq:theta-decomp}
\end{equation}
and rewrites the Wald statistic $T = \hat\theta/\widehat{\mathrm{SE}}$
in terms of two dimensionless channels,
\begin{equation}
T \approx \frac{\mu + \sigma Z^\ast}{\kappa\sigma} =
  \frac{1}{\kappa}\big(Z^\ast + \delta\big), \qquad
  \kappa = \frac{\overline{\widehat{\mathrm{SE}}}}{\sigma}, \qquad
  \delta = \frac{\mu}{\sigma}.
\label{eq:theta-wald}
\end{equation}
Realized Type~I error follows directly:
\begin{equation}
\alpha_{\mathrm{real}} = \Pr\big(|T| > t_{\nu,1-\alpha/2}\big) \approx
  \Pr\big(|Z^\ast + \delta| > \kappa\, t_{\nu,1-\alpha/2}\big),
\label{eq:theta-alpha}
\end{equation}
and exactly \emph{four} channels can move it away from nominal: a
standardized bias $\delta$, a scale-distortion $\kappa$ (model-based
SE too large or small relative to the true sampling SD), a
degrees-of-freedom channel $\nu$, and the shape of $Z^\ast$'s
distribution (departures from normality). Calibration is recovered
exactly when $\delta=0$, $\kappa=1$, $\nu=\infty$, and $Z^\ast\sim
N(0,1)$. With the bias and reference-distribution channels inert, this
reduces to a closed form in $\kappa$ alone,
\begin{equation}
\alpha_{\mathrm{real}} \approx 2\,\Phi\big(-\kappa\, t_{\nu,1-\alpha/2}\big)
  \;\approx\; 0.05 - 0.23\,(\kappa - 1) \ \text{near } \kappa=1,
\label{eq:theta-closed}
\end{equation}
with elasticity $d\log\alpha_{\mathrm{real}}/d\log\kappa \approx -4.6$
at $\kappa=1$: a small overstatement of the model-based standard error
produces a disproportionate collapse in realized Type~I error, which
is the compendium's general explanation for why an idealized OLS
analysis that ignores the participant random effect is anti-
conservative --- its scale-distortion factor is approximately
\begin{equation}
\kappa_{\mathrm{OLS}}^2 \approx \frac{\sigma_b^2 +
  \sigma_\varepsilon^2}{\sigma_\varepsilon^2},
\label{eq:kappa-ols}
\end{equation}
exceeding one by exactly the random-intercept variance it omits.

Paper 02 gives the compendium's operational estimator of $\kappa$
directly on the interaction coefficient, using the notation
$\hat\beta_{bm:X}$ for the fitted coefficient under carryover
specification $X$:
\begin{equation}
\kappa = \frac{\overline{\widehat{SE}}}{\mathrm{SD}(\hat\beta_{bm:X})}
  = \frac{\frac{1}{n_{\mathrm{sim}}}\sum_{r=1}^{n_{\mathrm{sim}}}
  \widehat{SE}_r}{\sqrt{\frac{1}{n_{\mathrm{sim}}-1}
  \sum_{r=1}^{n_{\mathrm{sim}}} (\hat\beta_{bm:X,r} -
  \bar{\hat\beta}_{bm:X})^2}},
\label{eq:kappa-estimator}
\end{equation}
and the calibrated true value against which bias is computed
throughout the carryover-specification manuscript,
\begin{equation}
\theta_{\mathrm{true}} = -c_{bm}\,\sigma_{BR}.
\label{eq:theta-true}
\end{equation}

\subsection{Cluster-robust sandwich variance and the CR2 correction}
\label{sec:sandwich}

The naive (Liang--Zeger) sandwich estimator, clustering on participant,
\begin{equation}
\widehat{\mathrm{Var}}_{\mathrm{sandwich}}(\hat\beta) =
  \Big(\textstyle\sum_i X_i' \hat V_i^{-1} X_i\Big)^{-1}
  \Big(\textstyle\sum_i X_i' \hat V_i^{-1} \hat e_i \hat e_i'
  \hat V_i^{-1} X_i\Big)
  \Big(\textstyle\sum_i X_i' \hat V_i^{-1} X_i\Big)^{-1},
\label{eq:sandwich}
\end{equation}
is downward-biased at moderate cluster counts. Paper 10 gives the
general CR family in the GLS-transformed metric,
\begin{equation}
\hat V_{\mathrm{CR}} = (\tilde X'\tilde X)^{-1}
  \Big(\textstyle\sum_i \tilde X_i' A_i \tilde e_i \tilde e_i' A_i
  \tilde X_i\Big) (\tilde X'\tilde X)^{-1},
\label{eq:cr-family}
\end{equation}
indexed by the adjustment matrix $A_i$: $A_i = I$ (CR0, unadjusted);
$A_i = (I-H_{ii})^{-1}$ (the jackknife form, equal to the
Mancl--DeRouen correction in the marginal/GEE setting); and
\begin{equation}
A_i = (I - H_{ii})^{-1/2} \qquad \text{(CR2)},
\label{eq:cr2}
\end{equation}
the symmetric square root that renders $\hat V_{\mathrm{CR}}$ exactly
unbiased when the working covariance is correctly specified --- the
cluster generalization of the HC2 estimator of MacKinnon and White.
Its Satterthwaite degrees of freedom are recovered by matching the
first two moments of the quadratic form's eigenvalues $\lambda_g$,
\begin{equation}
\nu = \frac{\big(\sum_g \lambda_g\big)^2}{\sum_g \lambda_g^2}.
\label{eq:cr2-df}
\end{equation}
Paper 02's null-cell diagnostics find $\kappa$ in the range
$1.05$--$1.10$ (empirical Type~I error $0.029$--$0.039$ against
nominal $0.05$) for the model-based standard errors, and $\kappa$
$1.07$--$1.26$ more broadly across specifications --- consistent, via
Equation~\eqref{eq:theta-closed}, with the mild conservatism the
compendium's Wald tests generally exhibit rather than
anti-conservatism.

%=============================================================
\section{The treatment main-effect track}
\label{sec:maineffect}
%=============================================================

Paper 04 targets a different estimand under a different, simpler DGP:
the treatment \emph{main} effect $\beta_D$, with \textbf{no biomarker
moderation channel at all}, and a linear (not Gompertz) data-
generating model.

\subsection{Data-generating model}

For the N-of-1 arm,
\begin{equation}
Y_{ij} = \beta_0 + \beta_D\, D_{bc,ij} + \beta_t\, t_{ij} + u_i +
  \varepsilon_{ij}, \qquad u_i \sim N(0,\sigma_u^2),
\label{eq:me-nof1}
\end{equation}
with $\boldsymbol\varepsilon_i$ following a continuous-time AR(1)
covariance $\mathrm{Cov}(\varepsilon_j,\varepsilon_k) = \sigma_w^2
\rho^{|t_j-t_k|}$; carryover is applied as an additive exponential
decay to the period immediately following a drug period, $\delta_c =
\beta_D \exp\{-\ln(2)\,L/t_{1/2}\}$ with $L$ the period length in
days. For the comparator parallel-group RCT arm, drawing patient
baselines from a shared pool,
\begin{equation}
Y_i = \mu_0 + u_i + \beta_D Z_i + \varepsilon_i, \qquad u_i \sim
  N(\mu_0, \sigma_\tau^2), \qquad Z_i \in \{0,1\}.
\label{eq:me-rct}
\end{equation}
The two arms share the same $N=35$ patient pool (participant-count-
matched, not observation-count-matched --- flagged by the paper itself
as a limitation), fit respectively by \code{nlme::lme(Sx \textasciitilde\ Dbc +
t\_wk, ..., correlation = corCAR1(...))} with Satterthwaite df, and by
ANCOVA, \code{lm(endpoint \textasciitilde\ trt + baseline)}.

\note{Internal inconsistency, resolved.} Paper 04 previously stated
$\sigma_w = 1.0$ in the AR(1) covariance definition of its DGP
(\S2.1) while separately citing $\sigma_w = 2.3$ as ``moderate
within-patient measurement noise'' in its optimal-design-conditions
discussion (\S6), with the two values never reconciled in the source
text. Inspection of the production driver
(\code{analysis/scripts/treatment-main-effect/production-run.R},
constant \code{SIGMA\_WITHIN <- 2.3}) confirms $\sigma_w = 2.3$ is
the value the simulation actually uses; \S2.1's $\sigma_w = 1.0$ was
the error, and has been corrected in the manuscript to match both
\S6 and the code.

\subsection{Classical closed-form power theory (background, not the
paper's own simulation)}

Paper 04's literature-review section reproduces the closed-form
power apparatus this compendium otherwise replaces with simulation,
due to Senn (2018) and Zucker et al.:
\begin{align}
d_{ij} &= \delta_i + \varepsilon_{ij}, \qquad \delta_i \sim
  N(\Delta,\tau^2), \qquad \varepsilon_{ij}\sim N(0, 2\sigma^2)
  \label{eq:senn-model} \\
\mathrm{Var}(\hat\Delta) &= \frac{\tau^2 + 2\sigma^2/k}{n} \qquad
  \text{(random-effects, patient-mean analysis)}
  \label{eq:senn-var-re} \\
\lambda_{\mathrm{RE}} &= \frac{\Delta}{\sqrt{(\tau^2+2\sigma^2/k)/n}},
  \quad df = n-1
  \label{eq:senn-ncp-re} \\
\lambda_{\mathrm{FE}} &= \frac{\Delta}{\sqrt{2\sigma^2/(nk)}}, \quad
  df = n(k-1) \qquad \text{(fixed-effects, pools within-patient SS)}
  \label{eq:senn-ncp-fe}
\end{align}
for $n$ patients and $k$ paired treatment cycles each, with $\tau^2$
the between-patient (treatment-by-patient interaction) variance and
$2\sigma^2$ the within-patient, within-cycle variance of paired
differences. The paper notes standard software assuming $n-1$
denominator df underestimates power here whenever $k>2$, because the
fixed-effects analysis correctly pools within-patient degrees of
freedom that the random-effects analysis discards. Zucker et al.'s
equivalent precision formula,
\begin{equation}
\mathrm{Var}(\hat\Delta) = \frac{\tau^2 + \sigma_w^2/k}{n},
\label{eq:zucker-var}
\end{equation}
makes explicit that the marginal benefit of additional cycles $k$
diminishes once $\sigma_w^2/k \ll \tau^2$: past that point, adding
cycles to existing patients buys far less precision than adding new
patients, because the between-patient variance floor dominates.

Paper 04's own simulation (2{,}000 replicates per cell, effect-size
grid $\beta_D\in\{0,-0.5,-1.0,-2.0\}$ nightmares/week, carryover
half-life sweep $t_{1/2}\in\{0.5,1.0,3.0,5.0,7.0\}$ days) does not
invoke Equations~\eqref{eq:senn-var-re}--\eqref{eq:zucker-var}
directly; it estimates power and Type~I error empirically, at the
MCSE rates given in \S\ref{sec:performance} below.

%=============================================================
\section{Shared performance measures}
\label{sec:performance}
%=============================================================

Every manuscript in the compendium reports the same three Monte Carlo
performance quantities, following the ADEMP framework of Morris,
White, and Crowther (2019):

\begin{align}
\hat\pi &= \frac{1}{n_{\mathrm{reps}}}\sum_{r=1}^{n_{\mathrm{reps}}}
  \mathbf 1[p_r < \alpha] \qquad \text{(empirical power / Type~I error)}
  \label{eq:perf-power} \\
\mathrm{MCSE}(\hat\pi) &= \sqrt{\hat\pi(1-\hat\pi)/n_{\mathrm{reps}}}
  \label{eq:perf-mcse-power} \\
\mathrm{Bias} &= \frac{1}{n_{\mathrm{reps}}}\sum_r \hat\theta_r -
  \theta_{\mathrm{true}}, \qquad \mathrm{MCSE(Bias)} =
  \frac{s_{\hat\theta}}{\sqrt{n_{\mathrm{reps}}}}
  \label{eq:perf-bias} \\
\mathrm{MCSE}(\widehat{SE}) &= \frac{s_{\hat\theta}}
  {\sqrt{2(n_{\mathrm{reps}}-1)}}
  \label{eq:perf-mcse-se}
\end{align}
with $s_{\hat\theta}$ the empirical standard deviation of the
replicate-level point estimates. A minimum replicate count of
$n_{\mathrm{reps}}\ge625$ is the compendium's recurring target,
derived from requiring $\mathrm{MCSE}(\hat\pi) < 0.02$ at the
worst-case $\hat\pi=0.5$; most production runs use $n_{\mathrm{reps}}
\in \{500, 1000\}$ for power/bias cells and up to $5000$ for
null/Type-I cells, where the smaller true $\hat\pi$ near $\alpha$
tightens the MCSE requirement at fixed sample size.

For a difference-of-differences comparison across two conditions
(e.g. the carryover-loss architecture contrast of \S\ref{sec:architectures}),
Paper 01 gives
\begin{equation}
\mathrm{SE}\big(\Delta_B - \Delta_A\big) = \sqrt{
  \frac{\hat\pi_{B,0}(1-\hat\pi_{B,0})}{n} +
  \frac{\hat\pi_{B,1}(1-\hat\pi_{B,1})}{n} +
  \frac{\hat\pi_{A,0}(1-\hat\pi_{A,0})}{n} +
  \frac{\hat\pi_{A,1}(1-\hat\pi_{A,1})}{n}}, \qquad \Delta_X =
  \hat\pi_{X,0} - \hat\pi_{X,1},
\label{eq:perf-dod}
\end{equation}
the standard error of the within-architecture carryover-loss gap,
used throughout the compendium to test whether two designs',
architectures', or specifications' power differences are
distinguishable from Monte Carlo noise.

%=============================================================
\section*{Summary: which manuscript supplies which piece}
\addcontentsline{toc}{section}{Summary: which manuscript supplies which piece}
%=============================================================

\begin{longtable}{@{}p{2.3cm}p{9.5cm}p{2.3cm}@{}}
\toprule
Paper & Contribution to this compendium & Key equations \\
\midrule
\endhead
01 & Mean- vs. covariance-moderation architectures; the central
  carryover-divergence result & \eqref{eq:archA}--\eqref{eq:archB-offdrug} \\
02 & Carryover decay families; nine analysis specifications E1--E9;
  $\kappa$ calibration; sandwich/CR2 & \eqref{eq:decay-exp}--\eqref{eq:e9-ols}, \eqref{eq:kappa-estimator}--\eqref{eq:theta-true} \\
03 & Latent-class reading of covariance moderation; MVN moment-matching
  approximation & \eqref{eq:latent-class}--\eqref{eq:bic-gate} \\
04 & Treatment main effect; simpler linear DGP; classical Senn/Zucker
  power theory & \eqref{eq:me-nof1}--\eqref{eq:zucker-var} \\
05 & Simple vs. complex designs; effective-$N$ shrinkage under AR(1);
  Gompertz and carryover-recursion restatements & \eqref{eq:effn-shrinkage}, \eqref{eq:gompertz-05}, \eqref{eq:carryover-recursion} \\
06 & Partitioning the treatment effect into $BR$/$PB$/$TV$; lumped-
  slope bias identity; expectancy scaling & \eqref{eq:decomposition}--\eqref{eq:phase-augmented} \\
07 & Modified-Gompertz curve and alternative trajectory families;
  architecture-dependent shape sensitivity & \eqref{eq:gompertz-07}--\eqref{eq:gompertz-cov-independence} \\
08 & Closed-form RM-ANOVA derivation; the compendium's only fully
  analytic power result & \eqref{eq:rmanova-model}--\eqref{eq:rmanova-power} \\
09 & Informative-dropout hazard; MNAR argument; design-dependent
  ``happy accident'' selection effect & \eqref{eq:dropout-hazard} \\
10 & General four-channel ($\delta,\kappa,\nu$, shape) decomposition
  of realized Type~I error; CR family & \eqref{eq:theta-decomp}--\eqref{eq:cr2-df} \\
11 & Dual-channel architecture; additive covariance, super-additive
  power, unstated non-identifiability & \eqref{eq:dualchannel-shift}--\eqref{eq:dc-superadditivity} \\
\bottomrule
\end{longtable}

\end{document}
